\subsection{\texorpdfstring{映射 \(f, v\) 和 \(\Phi\)}{映射 f、v 和 Φ}}

  设 A 为一个环。在 \(\mathbf{A}^{\mathbf{N}}\) 上赋予乘积环结构。以 \(f_{\mathbf{A}}\) 或简称 \(f\) 表示自同态 \((a_{n})_{n\in\mathbf{N}}\mapsto(a_{n+1})_{n\in\mathbf{N}}\) 的 \(\mathbf{A}^{\mathbf{N}}\)。以 \(v_{\mathbf{A}}\) 或简称 \(v\) 表示 \(\mathbf{A}^{\mathbf{N}}\) 的底层加法群的自同态，它将 \((a_{n})_{n\in\mathbb{N}}\) 送到 \((0,p\cdot a_{0},p\cdot a_{1},...)\)。

  对于每个整数 \(m \geqslant 0\)，以 \(\Phi_{m}\) 表示从 \(\mathbf{A}^{\mathbf{N}}\) 到 A 的映射，它将 \(\boldsymbol{a} = (a_{n})_{n \in \mathbf{N}}\) 送到 \(\Phi_{m}(a_{0}, \ldots, a_{m})\)。以 \(\Phi_{\mathrm{A}}\) 或简称 \(\Phi\) 表示映射 \(\boldsymbol{a} \mapsto (\Phi_{n}(\boldsymbol{a}))_{n \in \mathbf{N}}\)，它从 \(\mathbf{A}^{\mathbf{N}}\) 到自身。

引理 2.--- 设 A 为带有自同态 σ 的环，对于每个 a ∈ A，σ(a) ≡ a\^{}p 模 p.A。设 n ≥ 1 为整数，a\_0, \ldots, a\_\{n-1\} 为 A 的元素。置 u\_i = Φ\_i(a\_0, \ldots, a\_i)，其中 0 ≤ i ≤ n - 1。设 u\_n 为 A 的一个元素。下列条件等价：

\begin{enumerate}
\def\labelenumi{\alph{enumi})}
\item
  存在 \(a_{n} \in A\)，使得 \(u_{n} = \Phi_{n}(a_{0}, \ldots, a_{n})\)。
\item
  \(\sigma(u_{n-1}) \equiv u_n \pmod{p^nA}\).
\end{enumerate}

对于 \(0 \leqslant i \leqslant n - 1\)，有 \(\sigma(a_i) \equiv a_i^p \pmod{pA}\)。将第 1 号的命题 1 应用于 \(\mathbf{J}_k = p^kA\)（\(k \in \mathbf{N}\)）且 \(m = 1\) 的情形，得到同余关系

\[
\Phi_ {n - 1} \left(\sigma \left(a _ {0}\right), \dots , \sigma \left(a _ {n - 1}\right)\right) \equiv \Phi_ {n - 1} \left(a _ {0} ^ {p}, \dots , a _ {n - 1} ^ {p}\right) \pmod{p ^ {n}A},\tag{7}
\]

即

\[
\sigma (u _ {n - 1}) \equiv \Phi_ {n - 1} (a _ {0} ^ {p},..., a _ {n - 1} ^ {p}) \pmod{p ^ {n}A}.\tag{8}
\]

现在根据公式 (2)，关系 \(u_n = \Phi_n(a_0, \ldots, a_n)\) 等价于

\[
u _ {n} = \Phi_ {n - 1} (a _ {0} ^ {p},..., a _ {n - 1} ^ {p}) + p ^ {n} a _ {n}.\tag{9}
\]

引理由此得证。

命题 2. --- 设 A 为一个环。

\begin{enumerate}
\def\labelenumi{\alph{enumi})}
\item
  若 \(p.1_{\mathbf{A}}\) 不是 \(\mathbf{A}\) 中的零因子，则映射 \(\Phi_{\mathbf{A}}\) 是单射。
\item
  若 \(p.1_{\mathbf{A}}\) 在 \(\mathbf{A}\) 中可逆，则映射 \(\Phi_{\mathbf{A}}\) 是双射。
\item
  若 \(\sigma\) 是环 A 的自同态，且满足 \(\sigma(a) \equiv a^p \bmod pA\) 对所有 \(a \in \mathbf{A}\) 成立，则 \(\mathbf{A}'\)，即 \(\Phi_{\mathbf{A}}\) 的像，是 \(\mathbf{A}^{\mathbf{N}}\) 的子环，并且在 \(f_{\mathbf{A}}\) 和 \(v_{\mathbf{A}}\) 下稳定。它是元素 \((u_n)_{n \in \mathbf{N}}\) 的集合，这些元素属于 \(\mathbf{A}^{\mathbf{N}}\) 且满足 \(\sigma(u_n) \equiv u_{n+1} \bmod p^{n+1}A\) 对所有 \(n \in \mathbf{N}\) 成立。
\end{enumerate}

若 \(\boldsymbol{a} = (a_n)_{n \in \mathbb{N}}\) 和 \(\boldsymbol{u} = (u_n)_{n \in \mathbb{N}}\) 是 \(A^{\mathbb{N}}\) 的元素，则根据公式 (2)，关系 \(\Phi_A(\boldsymbol{a}) = \boldsymbol{u}\) 等价于等式

\[
\left\{ \begin{array}{l} u _ {0} = a _ {0}, \\ u _ {n} = \Phi_ {n - 1} \left(a _ {0} ^ {p}, \dots , a _ {n - 1} ^ {p}\right) + p ^ {n} a _ {n} \quad \text {for all} \quad n \geqslant 1. \end{array} \right.\tag{10}
\]

  设 \(\boldsymbol{u} = (u_{n})_{n \in \mathbb{N}}\) 为 \(A^{\mathbb{N}}\) 的元素。当 \(p1_{A}\) 是 A 中的非零因子（分别地，当 \(p1_{A}\) 在 A 中可逆）时，A 中至多存在一个满足等式 (10) 的序列 \((a_{n})_{n \in \mathbb{N}}\)（分别地，恰有一个 A 中满足等式 (10) 的序列 \((a_{n})_{n \in \mathbb{N}}\)），从而得到 a) 和 b)。

 我们证明 \(c)\)。根据引理 2，\(\mathbf{A}'\)，即 \(\Phi_{\mathbf{A}}\) 的像，是元素 \(\boldsymbol{u} = (u_n)_{n \in \mathbb{N}}\) 的集合，这些元素属于 \(\mathbf{A}^{\mathbf{N}}\) 且满足 \(\sigma(u_n) \equiv u_{n+1} \bmod p^{n+1}A\) 对所有 \(n \in \mathbb{N}\) 成立。立即可见 \(\mathbf{A}'\) 是 \(\mathbf{A}^{\mathbf{N}}\) 的子环，在 \(f_{\mathbf{A}}\) 和 \(v_{\mathbf{A}}\) 下稳定。

注记。—设 \(\boldsymbol{a} = (a_n)_{n \in \mathbb{N}}\) 和 \(\boldsymbol{u} = (u_n)_{n \in \mathbb{N}}\) 为 \(A^{\mathbf{N}}\) 中满足 \(\boldsymbol{u} = \Phi_{\mathbf{A}}(\boldsymbol{a})\) 的元素，且 \(m\) 为整数 \(\geqslant 0\)。由 (10) 可推出下列断言：

若 \(u_n\) 在 \(0 \leqslant n \leqslant m\) 时属于 A 的子环 B，且对于每个 \(x \in A\)，关系 \(px \in B\) 蕴含 \(x \in B\)，则 \(a_n\) 在 \(0 \leqslant n \leqslant m\) 时属于 B。

若 A 带有 \(\mathbf{N}\)-分次，\(p1_{A}\) 是 A 中的非零因子，\(d \in \mathbf{N}\)，且 \(u_{n}\) 是次数为 \(dp^{n}\) 的齐次元素，对 \(0 \leqslant n \leqslant m\) 成立，则 \(a_{n}\) 是次数为 \(dp^{n}\) 的齐次元素，对 \(0 \leqslant n \leqslant m\) 成立。

